% Options for packages loaded elsewhere
\PassOptionsToPackage{unicode}{hyperref}
\PassOptionsToPackage{hyphens}{url}
%
\documentclass[
]{article}
\title{Clinical Studies}
\author{Michaelle Valenzuela}
\date{diciembre 2022}

\usepackage{amsmath,amssymb}
\usepackage{lmodern}
\usepackage{iftex}
\ifPDFTeX
  \usepackage[T1]{fontenc}
  \usepackage[utf8]{inputenc}
  \usepackage{textcomp} % provide euro and other symbols
\else % if luatex or xetex
  \usepackage{unicode-math}
  \defaultfontfeatures{Scale=MatchLowercase}
  \defaultfontfeatures[\rmfamily]{Ligatures=TeX,Scale=1}
\fi
% Use upquote if available, for straight quotes in verbatim environments
\IfFileExists{upquote.sty}{\usepackage{upquote}}{}
\IfFileExists{microtype.sty}{% use microtype if available
  \usepackage[]{microtype}
  \UseMicrotypeSet[protrusion]{basicmath} % disable protrusion for tt fonts
}{}
\makeatletter
\@ifundefined{KOMAClassName}{% if non-KOMA class
  \IfFileExists{parskip.sty}{%
    \usepackage{parskip}
  }{% else
    \setlength{\parindent}{0pt}
    \setlength{\parskip}{6pt plus 2pt minus 1pt}}
}{% if KOMA class
  \KOMAoptions{parskip=half}}
\makeatother
\usepackage{xcolor}
\IfFileExists{xurl.sty}{\usepackage{xurl}}{} % add URL line breaks if available
\IfFileExists{bookmark.sty}{\usepackage{bookmark}}{\usepackage{hyperref}}
\hypersetup{
  pdftitle={Clinical Studies},
  pdfauthor={Michaelle Valenzuela},
  hidelinks,
  pdfcreator={LaTeX via pandoc}}
\urlstyle{same} % disable monospaced font for URLs
\usepackage[margin=1in]{geometry}
\usepackage{graphicx}
\makeatletter
\def\maxwidth{\ifdim\Gin@nat@width>\linewidth\linewidth\else\Gin@nat@width\fi}
\def\maxheight{\ifdim\Gin@nat@height>\textheight\textheight\else\Gin@nat@height\fi}
\makeatother
% Scale images if necessary, so that they will not overflow the page
% margins by default, and it is still possible to overwrite the defaults
% using explicit options in \includegraphics[width, height, ...]{}
\setkeys{Gin}{width=\maxwidth,height=\maxheight,keepaspectratio}
% Set default figure placement to htbp
\makeatletter
\def\fps@figure{htbp}
\makeatother
\setlength{\emergencystretch}{3em} % prevent overfull lines
\providecommand{\tightlist}{%
  \setlength{\itemsep}{0pt}\setlength{\parskip}{0pt}}
\setcounter{secnumdepth}{-\maxdimen} % remove section numbering
\ifLuaTeX
  \usepackage{selnolig}  % disable illegal ligatures
\fi

\begin{document}
\maketitle

\includegraphics{C:/Users/Dell/Documents/PEC3_Desarrollo/clinical_studies/Imagenes/logo.PNG}

El conjunto de datos corresponde a los \textbf{estudios clínicos de
cáncer de mama} que se encuentran en estado activos, terminados y
completos, de la fuente de ClinialTrial.gov.

\textbf{NTC:}

El código de identificación único que se le da a cada estudio clínico al
registrarse en ClinicalTials.goc. El formato del código está constituido
por ``NTC'' seguido de un número de 8 dígitos. Ejemplo: NCT00000419.

\textbf{Número Título:}

El título oficial de un protocolo utilizado para identificar un estudio
clínico o un título breve escrito en un lenguaje destinado al público no
especializado.

\textbf{Estado:}

Indica el estado de contratación actual o el estado de acceso ampliado.

\textbf{Resultado:}

Un registro del estudio que incluye los resultados resumidos publicados
en la base de datos de resultados de ClinicalTrials.gov. La información
resumida de los resultados incluye el flujo de participantes, las
características iniciales, las medidas de resultado y los eventos
adversos.

\textbf{Resultado del estudio:}

Un registro del estudio que incluye los resultados resumidos publicados
en la base de datos de ClinicalTrials.gov. La información resumida de
los resultados incluye el flujo de participantes, las características
iniciales, las medidas de resultado y los eventos adversos.

\textbf{Condiciones:}

La enfermedad, trastorno, síndrome, enfermedad o lesión que se está
estudiando.

\textbf{Intervenciones:}

Un proceso o acción que es el foco de un estudio clínico. Las
intervenciones incluyen medicamentos, dispositivos médicos,
procedimientos, vacunas y otros productos que están en investigación o
que ya están disponible. Las Intervenciones también pueden incluir
enfoques no invasivos como la educación o la modificación de la dieta y
el ejercicio.

\textbf{Patrocinador / Colaborador:}

Una organización distinta del patrocinador que brinda apoyo para un
estudio clínico. Este apoyo puede incluir actividades relacionadas con
la financiación, el diseño, la implementación, el análisis de datos o la
presentación de informes.

\textbf{Sexo:}

Un tipo de criterio de elegibilidad que indica el sexo de las personas
que pueden participar en un estudio clínico (todos, mujeres, hombres).
El sexo es la clasificación de una persona como hombre o mujer basada en
distinciones biológicas. El sexo es distinto de la elegibilidad basada
en el género.

\textbf{Edad:}

Tipo de criterio de elegibilidad que indica la edad que debe tener una
persona para participar en un estudio clínico. Esto puede estar indicado
por una edad específica o los siguientes grupos de edad:

\begin{itemize}
\item
  Adulto mayor (65 +)
\item
  Adulto (18 - 64)
\item
  Niño (nacimiento - 17)
\end{itemize}

\textbf{Fases:}

La etapa de un ensayo clínico que estudia un medicamento o producto
biológico, según las definiciones desarrolladas por la Administración de
Drogas y Alimentos de los EE.UU. (FDA). La fase se basa en el objetivo
del estudio, el número de participantes y otras características. Hay
cinco fases: Fase 0 o inicial, Fase 1, Fase 2, Fase 3 y Fase 4. No
aplicable se utiliza para describir ensayos sin fases definidas por la
FDA, incluidos los ensayos de dispositivos o intervenciones
conductuales.

\textbf{Inscripciones:} El número de participantes en un estudio
clínico.

\textbf{Tipo de financiador:}

Describe la organización que proporciona financiación o apoyo para un
estudio clínico. Este apoyo puede incluir actividades relacionadas con
la financiación, el diseño la implementación, el análisis de datos o la
presentación de informes. Las organizaciones que figuran como
patrocinadores y colaboradores de un estudio se consideran financiadores
del estudio. ClinicalTrials.gov cuenta con cuatro tipos de
financiadores:

\begin{itemize}
\item
  Institutos Nacionales de Salud de EE.UU
\item
  Otras agencias federales de EE.UU. (Administración de Alimentos y
  Medicamentos, los Centros para el Control y la Prevención de
  Enfermedades o el Departamento de Asuntos de Veteranos de EE.UU)
\item
  Industria (empresas farmacéuticas y de dispositivos)
\item
  Todos los demás (incluidos individuos, universidades y organizaciones
  comunitarias)
\end{itemize}

\textbf{Tipo de estudio:}

Describe la naturaleza de un estudio clínico. Los tipos de estudios
incluyen los ensayos clínicos, estudios observacionales y acceso
ampliado.

\textbf{Fecha de inicio:} La fecha real en la que se inscribió al primer
participante en un estudio clínico.

\textbf{Fecha de finalización primaria:}

La fecha en la que el último participante en un estudio clínico fue
examinado o recibió una intervención para recopilar datos finales para
las medidas de resultado primaria. Si el estudio clínico finalizó de
acuerdo con el protocolo o se dio por terminado no afecta esta fecha.
Para los estudios clínicos con más de una medida de resultado primaria
con diferentes fechas de finalización, este término se refiere a la
fecha en la que se completa la recopilación de datos para todas las
medidas de resultado primarias.

\textbf{Fecha primera publicación:}

La fecha en la que la información de los resultados resumidos estuvo
disponible por primera vez en ClinicalTrials.gov después de que concluyó
la revisión del control de calidad (QC) de la Biblioteca Nacional de
Medicina (NLM). Por lo general, hay una demora entre la fecha en que el
patrocinador del estudio o el investigador envía por primera vez la
información resumida de los resultados y la fecha de publicación de los
resultados.

\textbf{Última actualización publicada:}

La fecha más reciente en la que se pusieron a disposición los cambios en
el registro de un estudio en ClinicalTrials.gov. Puede haber una demora
entre el momento en que el patrocinador o investigador del estudio envió
los cambios a ClinicalTrials.gov y la fecha de publicación de la última
actualización.

\textbf{Localización:}

Países en los que se encuentran las instalaciones de investigación para
un estudio. Un país aparece en la lista solo una vez, incluso si hay más
de una instalación en el país. La lista incluye todos los países a la
fecha de la última actualización presentada; cualquier país para el que
se eliminaron todas las instalaciones del registro del estudio que se
enumera en países de ubicación eliminados.

\textbf{URL:} Enlace de publicación de los estudios clínicos.

\end{document}
